\begin{kbvderiv}
\vspace{4pt plus 3pt minus 1pt}
\noindent{\renewcommand{\arraystretch}{1.24}\setlength{\tabcolsep}{5pt}
\begin{tabularx}{\textwidth}{X X X}
ಷಟ್‌ + ಚಕ್ರ & ವಾಕ್‌ + ದೇವಿ & ಪೂರ್ವ + ಉತ್ತರ \\
ದುರ್ಗ + ಅಷ್ಟಮಿ & ಅತಿ + ಆಸೆ & ಯಶಸ್‌ + ಚಂದ್ರಿಕೆ \\
\end{tabularx}}
\vspace{4pt plus 3pt minus 1pt}
\end{kbvderiv}

\vspace{2.5pt plus 3pt minus 0.5pt}\noindent\hspace*{1.2em}\kbvnum{(11)}\hspace{0.55em}ಇವುಗಳನ್ನು ಬಿಡಿಸಿ ಬರೆದು ಸಂಧಿಯನ್ನು ತಿಳಿಸಿ. ಧನಾರ್ಜನೆ, ಮಹೇಶ್ವರ, ಸಿದ್ಧೌಷಧಿ, ಪ್ರತ್ಯುತ್ತರ... ಜಗದ್ರಕ್ಷಕ, ಷಡ್ದರ್ಶನ, ಚಿದಾನಂದ, ಮನಶ್ಶಾಂತಿ ಕಠಶೂಕ್ತಿ, ವೀರೇಂದ್ರ, ಮಹೌನ್ನತ್ಯ, ಷಡ್ಭುಜ - ಸಚ್ಚರಿತ್ರ, ಅಬ್ಧಿ, ದಿಗಂತ, ಅತ್ಯವಸರ\par

\vspace{2.5pt plus 3pt minus 0.5pt}\noindent\hspace*{1.2em}\kbvnum{(12)}\hspace{0.55em}. ಇವುಗಳನ್ನು ಬಿಡಿಸಿ ಬರೆಯಿರಿ ; ಬ್ರಹ್ಮರ್ಷಿ, ವೀರ್ಯಾಣು, ಅತುಳೈಶ್ವರ್ಯ, ದೇವಾಲಯ ಷಡ್ದಳ ಜನೈಕ್ಯತೆ, ಮನಶ್ಚಾಪಲ್ಯ, . ವಿಜಯೋತ್ಸವ\par

\kbvchapter{ಪಾಠ ೮}{ನಾಮಪದದ ವಿಭೇದಗಳು}

\indent “ರವಿಯು ತೋಟಕ್ಕೆ ನುಗ್ಗಿದ ಎಮ್ಮೆಯನ್ನು ಕೋಲಿನಿಂದ ಥಳಿಸಿದನು.”\par

\indent —ಈ ವಾಕ್ಯದಲ್ಲಿ ರವಿಯು, ತೋಟಕ್ಕೆ, ಎಮ್ಮೆಯನ್ನು, ಕೋಲಿನಿಂದ ಎಂಬ ಪದಗಳು ನಾಮಪದಗಳಾಗಿರುವುವು... ಈ ನಾಮಪದಗಳಲ್ಲಿ, ರವಿ, ತೋಟ ಎಮ್ಮೆ , ಕೋಲು. ಇವು ನಾಮಪ್ರಕೃತಿಗಳು ' ಮನುಷ್ಯ, ಸ್ಥಳ ಪ್ರಾಣಿ, ಪದಾರ್ಥ ಮುಂತಾದವುಗಳ ಹೆಸರುಗಳನ್ನು ತಿಳಿಸುವ ಈ ನಾಮಪ್ರಕೃತಿಗಳಿಗೆ “ನಾಮವಾಚಕ ಪ್ರಕೃತಿ” ಗಳೆಂದೂ ಹೆಸರು... ನಾಮವಾಚಕ ಪ್ರಕೃತಿಗಳಲ್ಲಿ ಹಲವು ಪ್ರಕಾರಗಳುಂಟು ; ಅವುಗಳಲ್ಲಿ ಮುಖ್ಯವಾದವುಗಳೆಂದರೆ—\par

\vspace{4pt plus 3pt minus 1pt}
\begin{center}
{\renewcommand{\arraystretch}{1.22}\setlength{\tabcolsep}{5pt}
\begin{tabular}{r l c l}
(೧) & ವಸ್ತುವಾಚಕಗಳು & ಉದಾ : & ಮರ, ಗೌರಿ \\
(೨) & ಗುಣವಾಚಕಗಳು & ಉದಾ : & `: : ದೊಡ್ಡ, ಹಳೆಯ \\
(೩) & ಸಂಖ್ಯಾವಾಚಕಗಳು & ಉದಾ : & ಒಂದು ಹತ್ತು \\
\end{tabular}}
\end{center}
\vspace{4pt plus 3pt minus 1pt}

\vspace{2.5pt plus 3pt minus 0.5pt}\noindent\hspace*{1.2em}\kbvnum{(೪)}\hspace{0.55em}. ಸಂಖ್ಯೇಯವಾಚಕಗಳು. ಉದಾ : ಐವರು, ಹತ್ತನೆಯ\par

\vspace{4pt plus 3pt minus 1pt}
\begin{center}
{\renewcommand{\arraystretch}{1.22}\setlength{\tabcolsep}{5pt}
\begin{tabular}{r l c l}
(೫) & ಪರಿಮಾಣವಾಚಕಗಳು & ಉದಾ : & ಅಷ್ಟು, ಇಷ್ಟು \\
(೬) & ಪ್ರಕಾರವಾಚಕಗಳು & ಉದಾ : & ಅಂತಹ, ಇಂತಹ \\
(೭) & ಸರ್ವನಾಮಗಳು & ಉದಾ : & ಅವನು ಇವನು \\
\end{tabular}}
\end{center}
\vspace{4pt plus 3pt minus 1pt}
